\section{Minimality}\label{sec:minimality}

Minimality asks which axioms of the six figure-defined presentations are genuinely necessary for
completeness, and the full and bounded statements in \cref{thm:minimality-all} are certified by
separating interpretations.  Throughout, fix a fragment \(\bullet\) with free PROP
\(\cat{P}_\bullet\) and finite rule set \(\QC_\bullet\).  The arity bounds in the theorem are dictated
by separator coverage: a bounded minimality claim is made only when the table contains a
non-\(\noseparator\) separator for every axiom in that truncation.  When this independence is read as
necessity for completeness, the Clifford\(+T\) and Clifford\(+CS\) cases are understood relative to the
bounded completeness results imported in \cref{sec:completeness}.

\subsection{Independence and minimality}

To test whether an axiom \(\rho\) is needed, remove it from the rule set and ask whether its two sides
are still equal in the quotient by the remaining rules.  We write
\(\QC_\bullet\setminus\{\rho\}\) for those remaining axioms, and
\(\QCirc_\bullet^{-\rho}\coloneqq \cat{P}_\bullet/(\QC_\bullet\setminus\{\rho\})\) for the reduced
presentation.

\begin{definition}\label{def:minimality}
Let \(\rho\) be an axiom \(C_1=C_2\) in \(\QC_\bullet\) (with \(C_1,C_2:n\to n\)).
We say that \(\rho\) is independent of the other axioms of the same fragment if
\(
  \QCirc_\bullet^{-\rho}\not\vdash C_1=C_2.
\)
We say that \(\QC_\bullet\) is minimal if every \(\rho\in\QC_\bullet\) is independent.
\end{definition}

Because the presentations are sound for their intended semantics, an axiom derivable from the
others cannot be necessary for completeness: removing it leaves the induced congruence unchanged.
Conversely, in any arity range where the reduced presentation were still complete, soundness would
force the removed equation to remain derivable.  Thus independence is the syntactic test used here,
and the separators in \cref{fig:minimality-all} instantiate it by occurrence detectors, counting
models, projective substitutions, and determinant phases.

Since derivations preserve arity by \cref{rem:arity-preservation}, an \(n\)-wire axiom can be tested
inside the subsystem generated by axioms of arity at most \(n\).  This gives the bounded notion used
when the separator table contains separators only in a finite arity range.

\begin{definition}\label{def:bounded-minimality}
Let \(k\in\mathbb{N}\).  Define
\(
  \QC_{\bullet,\le k}
   \coloneqq 
  \{ \rho\in\QC_\bullet \mid \rho \text{ has type } n\to n \text{ with } n\le k \}.
\)
We say that \(\QC_\bullet\) is minimal up to \(k\) wires if every \(\rho\in\QC_{\bullet,\le k}\) is independent
relative to \(\QC_{\bullet,\le k}\) (i.e.\ independent in the sense of \cref{def:minimality} after replacing
\(\QC_\bullet\) by \(\QC_{\bullet,\le k}\)).
\end{definition}

\subsection{Separation by alternative interpretations}

Our independence proofs are semantic: to show that an axiom \(\rho\) is not derivable from the others,
we construct a PROP-valued model that satisfies all remaining axioms but not \(\rho\).  The underlying
countermodel principle for equational logic is that derivable equations hold in every model of the
theory, so one model of the reduced theory that violates \(\rho\) rules out a derivation.  This
principle goes back at least to Birkhoff's correspondence between equational theories and varieties
of algebras \cite{Birkhoff1935} and is explicit in modern accounts of equational completeness
\cite[Thm.~14.19]{BurrisSankappanavar1981}.  In the quantum-circuit literature, the same proof
pattern is used in \cite{ClementDelormePerdrix2023}, which proves a minimal and complete
equational theory for general quantum circuits.

\begin{lemma}\label{lem:separation}
Let \(\rho:C_1=C_2\) be an axiom in \(\QC_\bullet\).  Assume there exist a PROP
\(\cat{P}_{\mathrm{alt}}\) and a PROP morphism \(F:\cat{P}_\bullet\to\cat{P}_{\mathrm{alt}}\) such that
\(F(L)=F(R)\) for every \((L=R)\in\QC_\bullet\setminus\{\rho\}\), but \(F(C_1)\neq F(C_2)\).  Then
\(\rho\) is independent.
\end{lemma}

\begin{proof}
Since \(F\) equalises every axiom in \(\QC_\bullet\setminus\{\rho\}\), it factors through the quotient
\(\QCirc_\bullet^{-\rho}\), yielding a PROP morphism
\(\overline{F}:\QCirc_\bullet^{-\rho}\to\cat{P}_{\mathrm{alt}}\).
If \(\QCirc_\bullet^{-\rho}\vdash C_1=C_2\) then applying \(\overline{F}\) gives \(F(C_1)=F(C_2)\), a contradiction.
\end{proof}

Because morphisms of the free PROP are already quotiented by coherence, once \(F\) is a PROP morphism
it automatically respects all structural equalities, so the separator check is purely
fragment-specific.

\subsection{Families of separating interpretations}

\Cref{fig:minimality-all} is the finite certificate for the independence proofs: each
non-\(\noseparator\) row supplies one separator to be checked against the other axioms in the
relevant arity range, while the \(\noseparator\) rows mark exactly the gaps that prevent full
minimality claims for qutrit Clifford, Clifford\(+T\), and Clifford\(+CS\).  The separators fall into
four families--counting models, occurrence detectors, projective substitutions, and
determinant-phase interpretations--so the proof remains uniform even though the fragments differ.
In displayed rule sets and tables, abbreviated gates stand for fixed morphisms of the free PROP, and
every separating interpretation is applied to the expanded morphisms.

\emph{Counting and occurrence detectors.}
Counting models record selected generator multiplicities or swap parity, while occurrence detectors
record whether a chosen generator appears; both kinds of separator land in endomorphism-only PROPs,
so there are no morphisms between different arities.

\begin{definition}\label{def:endM}
Let \((M,\oplus,0)\) be a commutative monoid and let \(\epsilon\in M\) satisfy
\(\epsilon\oplus\epsilon=0\).  Define \(\cat{End}_{\epsilon}(M)\) to be the PROP with objects
\(\mathbb{N}\) and hom-sets
\(
  \cat{End}_{\epsilon}(M)(n,n)\coloneqq M \; (n\in\mathbb{N}),
\)
\(
  \cat{End}_{\epsilon}(M)(n,m)\coloneqq \varnothing \; (n\neq m).
\)
Composition and tensor are given by \(\oplus\), identities by \(0\), and the basic symmetry
\(\sigma_{1,1}\) is interpreted as \(\epsilon\).
\end{definition}

\begin{remark}
The condition \(\epsilon\oplus\epsilon=0\) implies \(\sigma_{m,n}\circ\sigma_{n,m}=\id_{n+m}\).
All other PROP coherence axioms hold because every structural map is interpreted as an element of
the commutative monoid \(M\), and both \(\circ\) and \(\otimes\) are interpreted as \(\oplus\).
\end{remark}

\begin{definition}\label{def:count-has}
For \(m\ge 2\) define
\(
  \cat{Count}_{m}\coloneqq \cat{End}_{0}(\mathbb{Z}_m,+,0)
\)
and
\(
  \cat{Has}\coloneqq \cat{End}_{0}(\mathbb{B},\vee,0),
\)
with \(\mathbb{B}=\{0,1\}\).
When we want to record the parity of the structural swap, we also use
\(
  \cat{Count}^{\mathrm{swap}}_{2}\coloneqq \cat{End}_{1}(\mathbb{Z}_2,+,0),
\)
so that the basic symmetry \(\sigma_{1,1}\) is interpreted as \(1\in\mathbb{Z}_2\).
\end{definition}

\begin{definition}\label{def:counting-notation}
Let \(m\ge 2\).  Given a finite list \((g_1,\ldots,g_r)\) of generators in \(\Sigma_\bullet\), with
repetitions counted with multiplicity, write
\(
  \#\{g_1,\ldots,g_r\}_{[m]}:\cat{P}_\bullet\to\cat{Count}_m
\)
for the unique PROP morphism sending each generator \(g\) to the number of its occurrences in the
list modulo \(m\).
\end{definition}

When \(m=2\), we use the symbol \(\gSWAP\) inside the braces only as a mnemonic that swaps are counted:
\(
  \#\{g_1,\ldots,g_r,\gSWAP\}_{[2]}:\cat{P}_\bullet\to\cat{Count}^{\mathrm{swap}}_2.
\)
In particular, \(\#\{\gSWAP\}_{[2]}\) denotes the morphism that sends every non-structural generator
to \(0\) and counts each structural swap as \(1\).

\begin{definition}\label{def:has-functors}
For a generator \(g\in\Sigma_\bullet\), write
\(
  ?g:\cat{P}_\bullet\to\cat{Has}
\)
for the unique PROP morphism with \(?g(g)=1\) and \(?g(h)=0\) for every generator \(h\neq g\).
Equivalently, for any circuit \(C\), one has \(?g(C)=1\) if and only if \(g\) occurs at least once in
\(C\).
\end{definition}

\begin{example}\label{ex:clifford-t-h2-separator}
In the one-wire Clifford\(+T\) truncation, the detector \(?\gH\) separates the axiom
\eqref{t-H2} from the remaining rules.  The equations \eqref{t-w8} and \eqref{t-Tpow8} contain no
\(\gH\) on either side, while \eqref{t-E} contains an occurrence of \(\gH\) on both sides.  The same is
true of \eqref{t-TX} after expanding the displayed shortcut \(\gX\).  Thus \(?\gH\) sends each
remaining equation in \(\QC_{\mathit{CliffT},\le 1}\) to either \(0=0\) or \(1=1\).  On the removed axiom,
however, \(?\gH(\gHsquare)=1\) and \(?\gH(\id_1)=0\), so \eqref{t-H2} is independent by
\cref{lem:separation}.
\end{example}

\emph{Projective substitutions.}
These separators alter one generator while quotienting only by global phase, which separates rows
whose two sides have the same generator multiplicities but differ after a projective change of one
generator.

\begin{definition}\label{def:projective-quotient}
For \(d\in\{2,3\}\), let \(\cat{Qudit}_d^{\sim}\) be the projective quotient of
\(\cat{Qudit}_d\) from \cref{def:qubit-qutrit-props}: its endomorphism hom-set
is \(\cat{Qudit}_d(n,n)/\mathrm{U}(1)\), where \(\mathrm{U}(1)\) acts by global
phase, and its non-endomorphism hom-sets remain empty.  Write
\(\Pi_d:\cat{Qudit}_d\to\cat{Qudit}_d^{\sim}\) for the quotient PROP morphism.
\end{definition}

\begin{definition}\label{def:projective-interpretation}
For a fragment \(\bullet\), set \(d(\bullet)=3\) for the qutrit Clifford
fragment and \(d(\bullet)=2\) for the qubit fragments.  Its usual projective
interpretation on the free presentation is
\(\interp{\cdot}_{\bullet,\sim}\coloneqq
\Pi_{d(\bullet)}\circ\interp{\cdot}^{\mathrm{free}}_\bullet\), regarded as a
PROP morphism \(\cat{P}_\bullet\to\cat{Qudit}_{d(\bullet)}^{\sim}\), where
\(\interp{\cdot}^{\mathrm{free}}_\bullet:\cat{P}_\bullet\to
\cat{Qudit}_{d(\bullet)}\) is the standard unitary interpretation before
quotienting by the fragment equations.
\end{definition}

\begin{definition}\label{def:projective-subst}
Given a fragment \(\bullet\), a generator \(g:n\to n\), and a circuit \(C:n\to n\) of that fragment, write
\(\interp{\cdot}_{\bullet,\sim}^{g:=C}\) for the unique PROP morphism
\(\cat{P}_\bullet\to\cat{Qudit}_{d(\bullet)}^{\sim}\) sending \(g\) to the already-evaluated
class \(\interp{C}_{\bullet,\sim}\) and every generator \(h\neq g\) to
\(\interp{h}_{\bullet,\sim}\).  Displayed abbreviations in \(C\) are expanded before this
interpretation.
\end{definition}

Thus a table entry \(\psubst{g}{C}\) means that every generator is interpreted projectively as
usual except \(g\), whose value is replaced by the projective value of \(C\).

\emph{Determinant phases.}
These separators record a scaled determinant phase for the qubit rows, chosen so as to be strictly
monoidal on a bounded number of wires.

\begin{definition}\label{def:argdetk}
Fix a qubit fragment, let \(k\ge 2\), and write \(\arg\det(U)\in\mathbb{R}/2\pi\mathbb{Z}\) for the
phase of \(\det(U)\).  Set
\(
  \epsilon_k\coloneqq 2^{k-2}\pi \bmod 2\pi,
\)
so \(\epsilon_2=\pi\) and \(\epsilon_k=0\) for \(k>2\).

Define a PROP morphism
\(
  \arg\det_k:\cat{P}_\bullet\to \cat{End}_{\epsilon_k}(\mathbb{R}/2\pi\mathbb{Z},+,0)
\)
on generators \(g:n\to n\) by
\(
  (\arg\det_k)(g)\coloneqq 2^{k-n}\arg\det(\interp{g}) \bmod 2\pi,
\)
and extend it to all circuits using preservation of \(\circ\), \(\otimes\), and symmetries.
\end{definition}

\begin{remark}
The scaling factor \(2^{k-n}\) is chosen so that \(\arg\det_k\) is compatible with tensoring on
arities \(\le k\), using \(\det(U\otimes V)=\det(U)^{2^m}\det(V)^{2^n}\) for \(U\) on \(n\) qubits and \(V\)
on \(m\) qubits.  The choice of \(\epsilon_k\) makes the image of the structural swap
\(\sigma_{1,1}\) coincide with its determinant phase after scaling: it contributes \(\pi\) when \(k=2\)
and becomes \(0\) for \(k>2\).
For equations on more than \(k\) wires, \(\arg\det_k\) is still used only as the generator-defined
PROP morphism above; the row-wise separator check is a check in
\(\cat{End}_{\epsilon_k}(\mathbb{R}/2\pi\mathbb{Z},+,0)\), not an assertion that the displayed
value is the scaled determinant of the whole higher-arity unitary.
\end{remark}

For each non-\(\noseparator\) row of \cref{fig:minimality-all}, the relation named on the left is
removed and the separator on the right equalises the other axioms in the relevant arity range while
distinguishing the removed relation.  The table records the separator rather than every value in the
finite check.  Detector and counting rows use the relevant commutative monoid,
projective substitution rows use the projective quotient from \cref{def:projective-interpretation},
and determinant rows use the scaled determinant phase from \cref{def:argdetk}.  When a higher-arity
separator is used in \cref{thm:minimality-all}, its check is against all axioms in the relevant arity
truncation, including same-arity rows marked \(\noseparator\) in \cref{fig:minimality-all}.

%%%%%%%%%%%%%%%%%%%%%%%%%%%%%%%%%%%%%%%%%%%%%%%%%%%%%
%%%%%%%%%%%%%%%%%%%%%%%%%%%%%%%%%%%%%%%%%%%%%%%%%%%%%
\begin{figure}[ht]
  \centering
  \begin{subfigure}[t]{0.31\textwidth}
    \centering
    \caption{Clifford}
    \begin{tabular}{|c|c|}
      \hline
      \textbf{Rel.} & \textbf{Separator} \\
      \hline
      \eqref{w8}   & \(?\gw\)   \\
      \eqref{H2}   & \(?\gH\)   \\
      \eqref{S4}   & \(?\gS\)   \\
      \eqref{E}    & \(\#\{\gH\}_{[2]}\) \\
      \eqref{Cs}   & \(?\gCNOT\)\\
      \eqref{B}    & \(\#\{\gSWAP\}_{[2]}\) \\
      \eqref{CZ}   & \(\#\{\gCNOT,\!\gSWAP\}_{[2]}\) \\
      \eqref{Inew} & \(\arg\det_{2}\) \\
      \hline
    \end{tabular}
  \end{subfigure}
  \begin{subfigure}[t]{0.31\textwidth}
    \centering
    \caption{Real Clifford}
    \begin{tabular}{|c|c|}
      \hline
      \textbf{Rel.}      & \textbf{Separator} \\
      \hline
      \eqref{realminus2}     & \(?\gminus\) \\
      \eqref{realH2}         & \(?\gH\) \\
      \eqref{realZ2}         & \(?\gZ\) \\
      \eqref{realF}          & \(\#\{\gminus\}_{[2]}\) \\
      \eqref{realCX2}        & \(?\gCNOT\) \\
      \eqref{realB}          & \(\#\{\gSWAP\}_{[2]}\) \\
      \eqref{realwC}         & \(\#\{\gZ\}_{[2]}\) \\
      \eqref{realCF}         & \(\psubst{\gZ}{\gI}\) \\
      \eqref{realXC}         & \(\psubst{\gH}{\gI}\) \\
      \eqref{realInew}       & \(\arg\det_{2}\) \\
      \hline
    \end{tabular}
  \end{subfigure}
  \begin{subfigure}[t]{0.31\textwidth}
    \centering
    \caption{Qutrit Clifford}
    \begin{tabular}{|c|c|}
      \hline
      \textbf{Rel.}            & \textbf{Separator} \\
      \hline
      \eqref{qt-omegapow12}        & \(?\gww\) \\
      \eqref{qt-hpow4}             & \(?\gH\) \\
      \eqref{qt-spow3}             & \(?\gS\) \\
      \eqref{qt-shpow3}            & \(\#\{\gH\}_{[2]}\) \\
      \eqref{qt-ssprime}           & \(\psubst{\gS}{\gSX}\)\footnote{See \cref{qt-prop:SSpnecessity}.} \\
      \eqref{qt-cnotremove}        & \(?\gCNOT\) \\
      \eqref{qt-swapdecomp}        & \(\#\{\gSWAP\}_{[2]}\) \\
      \eqref{qt-czdecomp}          & \(\#\{\gS\}_{[3]}\) \\
      \eqref{qt-kcnot}             & \(\#\{\gCNOT\}_{[2]}\) \\
      \eqref{qt-I}             & \noseparator \\
      \hline
    \end{tabular}
  \end{subfigure}
  \begin{subfigure}[t]{0.35\textwidth}
    \centering
    \caption{Clifford+T}
    \begin{tabular}{|c|c|}
      \hline
      \textbf{Rel.}     & \textbf{Separator} \\
      \hline
      \eqref{t-w8}         & \(?\gw\) \\
      \eqref{t-H2}         & \(?\gH\) \\
      \eqref{t-Tpow8}      & \(?\gT\) \\
      \eqref{t-E}          & \(\#\{\gH\}_{[2]}\) \\
      \eqref{t-TX}         & \(\#\{\gH,\!\gT,\!\gw\}_{[2]}\) \\
      \eqref{t-Ct}         & \(?\gCNOT\) \\
      \eqref{t-B}          & \(\#\{\gSWAP\}_{[2]}\) \\
      \eqref{t-CZ}         & \(\#\{\gCNOT,\!\gSWAP\}_{[2]}\) \\
      \eqref{t-CsCh}         & \noseparator \\
      \eqref{t-HTHT}         & \noseparator \\
      \eqref{t-C20}         & \noseparator \\
      \hline
    \end{tabular}
  \end{subfigure}
  \begin{subfigure}[t]{0.31\textwidth}
    \centering
    \caption{Clifford+CS}
    \begin{tabular}{|c|c|}
      \hline
      \textbf{Rel.}     & \textbf{Separator} \\
      \hline
      \eqref{cs-w8}         & \(?\gw\) \\
      \eqref{cs-H2}         & \(?\gH\) \\
      \eqref{cs-S4}         & \(?\gS\) \\
      \eqref{cs-E}          & \(\#\{\gH\}_{[2]}\) \\
      \eqref{cs-C}          & \(?\gCS\) \\
      \eqref{cs-B}          & \(\#\{\gSWAP\}_{[2]}\) \\
      \eqref{cs-XCS}        & \(\#\{\gS,\!\gH\}_{[2]}\) \\
      \eqref{cs-CSrev}      & \(\psubst{\gCS}{\gCSz}\) \\
      \eqref{cs-SHCHC}      & \(\psubst{\gCS}{\gCSzz}\) \\
      \eqref{cs-I}          & \(\arg\det_{2}\) \\
      \eqref{cs-U}          & \noseparator \\
      \hline
    \end{tabular}
  \end{subfigure}
  \begin{subfigure}[t]{0.32\textwidth}
    \centering
    \caption{CNOT-dihedral}
    \begin{tabular}{|c|c|}
      \hline
      \textbf{Rel.} & \textbf{Separator} \\
      \hline
      \eqref{cnot-R10} & \(?\gw\) \\
      \eqref{cnot-R7}  & \(?\gT\) \\
      \eqref{cnot-R1}  & \(?\gX\) \\
      \eqref{cnot-R11} & \(\#\{\gw\}_{[2]}\) \\
      \eqref{new-R3}   & \(\#\{\gX\}_{[2]}\) \\
      \eqref{cnot-R12} & \(?\gCNOT\) \\
      \eqref{new-R5}   & \(\#\{\gSWAP\}_{[2]}\) \\
      \eqref{new-R8}   & \(\#\{\gT,\!\gw,\!\gw\}_{[8]}\) \\
      \eqref{cnot-R6}  & \(\#\{\gCNOT,\!\gSWAP\}_{[2]}\) \\
      \eqref{cnot-R9}  & \(\#\{\gT,\!\gw,\!\gw\}_{[4]}\) \\
      \eqref{cnot-R13} & \(\arg\det_{3}\) \\
      \hline
    \end{tabular}
  \end{subfigure}
  \caption{Separators witnessing independence, indexed by the equation labels from \cref{sec:circuits-relations}; entries marked \(\noseparator\) are the cases for which no separator is currently provided.}\label{fig:minimality-all}
\end{figure}

\subsection{Minimality statements}

The \(\noseparator\) rows lie outside the bounded ranges claimed in \cref{thm:minimality-all}.  When
the table contains a separator beyond such a range, the corresponding higher-arity axiom is stated
separately.

\begin{theorem}\label{thm:minimality-all}
The indicated presented circuit fragments have minimal rule sets as listed:
\begin{enumerate}
\item \(\QCclifford\) is minimal for \(\cat{P}_{\mathit{Cliff}}\)-circuits (all arities).
\item \(\QCrealclifford\) is minimal for \(\cat{P}_{\mathit{RCliff}}\)-circuits (all arities).
\item \(\QCqutritclifford\) is minimal for \(\cat{P}_{\mathit{Cliff3}}\)-circuits up to \(2\) qutrit wires.\footnote{Full minimality would follow from a separator for the \(3\)-qutrit axiom \cref{qt-I}.}
\item \(\QCcliffordplust\) is minimal for \(\cat{P}_{\mathit{CliffT}}\)-circuits up to \(1\) qubit wire.
\item \(\QCcliffordpluscs\) is minimal for \(\cat{P}_{\mathit{CliffCS}}\)-circuits up to \(2\) qubit wires.
\item \(\QCdihedral\) is minimal for \(\cat{P}_{\mathit{CNOTdihe}}\)-circuits (all arities).
\end{enumerate}

For the bounded fragments, the constructed higher-arity separators give:
\begin{itemize}
\item In \(\QCcliffordplust\), the \(2\)-qubit axioms \eqref{t-Ct}, \eqref{t-B} and \eqref{t-CZ}
      are independent.
\item In \(\QCcliffordpluscs\), the \(3\)-qubit axiom \eqref{cs-I} is independent.
\end{itemize}
\end{theorem}

\begin{proof}
For each axiom \(\rho\) covered by the statement, the corresponding non-\(\noseparator\) row of
\cref{fig:minimality-all} gives a separator \(F_\rho\) that equalises the
remaining axioms of the relevant truncated system while distinguishing the two sides of \(\rho\).
For the additional higher-arity statements, the relevant truncated system is the one at the arity of
the displayed axiom.  Apply \cref{lem:separation} and
\cref{def:bounded-minimality}.
\end{proof}
